\documentclass[11pt]{article}
\usepackage{icmlko}
\begin{document}

\title{여덟 중 하나가 아니라 열여덟 중 하나 --- 무리 안 판의 절반은
자의 성질이었다}
\author{월드모델 실험실}
\date{노트 485 · 2026년 8월 3일}
\maketitle

\begin{abstract}
노트 484 가 경고를 하나 남겼다 --- 애니에서 격차의 $62\%$ 가 \textbf{라벨
압축}이었으니, 노트 470 의 \textbf{무리 안 판}도 같은 병일 수 있다. 그때
대조는 \emph{같은 크기 무작위}였는데 그것은 \textbf{전체 라벨 퍼짐을
보존한다}. 갈래 무리는 크기만 작은 게 아니라 라벨도 좁다.

갈래 무리마다 \textbf{라벨 SD$\cdot$평균을 맞춘 같은 크기 부분집합}을
도메인 전체에서 뽑아 대조로 놓았다.

\begin{center}
\begin{tabular}{lrl}
\hline
 & 가중평균 rho & 몫 \\
\hline
전체 판 & $0.4963$ & \\
무작위 같은 크기 & $0.4895$ & 크기 $0.0068$ \\
\textbf{퍼짐 맞춤} & $\mathbf{0.4643}$ & \textbf{압축 $0.0252$} \\
무리 안 판 & $0.4377$ & 남은 것 $0.0266$ \\
\hline
\end{tabular}
\end{center}

\textbf{낙폭 $0.0586$ 의 $55\%$ 가 라벨 압축이다.} 노트 470 이 ``판의
여덟 중 하나가 무리 사이 순위''라 적었는데, 실제로 무리 사이에 기대는
것은 $\mathbf{0.0266}$ --- \textbf{열여덟 중 하나}다.
\end{abstract}

\begin{figure}[t]
\centering
\includegraphics[width=\textwidth]{figs/f1.pdf}
\caption{\textbf{가} --- 낙폭의 세 조각. \textbf{나} --- 퍼짐을 맞춘 뒤
도메인마다 남는 것.}
\end{figure}

\section{남은 절반은 어디 있나}

퍼짐을 맞춘 뒤에도 남는 $0.0266$ 은 \textbf{고르게 퍼져 있지 않다}.

\begin{center}
\begin{tabular}{lrrr}
\hline
도메인 & 퍼짐 맞춤 & 무리 안 & 남은 것 \\
\hline
\textbf{펀딩} & $+0.3207$ & $+0.2301$ & $\mathbf{-0.0906}$ \\
애니 & $+0.5308$ & $+0.4880$ & $-0.0428$ \\
모바일 & $+0.5975$ & $+0.5548$ & $-0.0427$ \\
세계애니 & $+0.5255$ & $+0.5450$ & $+0.0195$ \\
웹툰 & $+0.4634$ & $+0.4590$ & $-0.0044$ \\
\textbf{만화} & $+0.3064$ & $+0.3613$ & $\mathbf{+0.0549}$ \\
\hline
\end{tabular}
\end{center}

\textbf{펀딩이 남은 것의 대부분을 낸다.} 노트 471 이 펀딩 점수의
$83\%$ 가 \texttt{category} 조회표라고 했다 --- 무리를 빼면 남는 게
없는 것이 당연하다. \emph{두 노트가 같은 사실을 두 방향에서 짚는다.}

\textbf{만화는 반대다}($+0.0549$). 라벨 퍼짐을 맞춘 무리보다 \emph{실제
갈래 무리}에서 더 잘한다 --- 만화의 축은 갈래 \emph{안}에서 더 잘 듣는다.
노트 470 에서도 만화만 무리 안이 더 높았는데($+0.0170$), 퍼짐을 맞추니
그 값이 세 배가 된다.

\section*{노트 470 을 고쳐 적는다}

\begin{quote}
\textbf{정정} --- 노트 470 의 ``판 점수의 $12\%$ 가 무리 사이 순위''는
\textbf{절반이 자의 성질}이었다. 무리를 나누면 라벨이 좁아지고, 좁은
라벨은 순위로 덜 갈린다. \textbf{실제로 무리 사이에 기대는 몫은
$5.4\%$}이며, 그것도 \textbf{펀딩 한 도메인이 대부분}이고 \textbf{만화는
반대 방향}이다.
\end{quote}

\textbf{자 자체는 남긴다.} 무리 안 판은 노트 473 에서 \texttt{it} $400$
을 잡아냈고 그 값은 이 정정으로 안 사라진다 --- 설정끼리 견줄 때는
양쪽의 라벨 퍼짐이 같기 때문이다. \textbf{바뀌는 것은 그 절대값을
``실력 부족''으로 읽으면 안 된다는 것}이다.

\section*{조항 --- 자를 만들 때 그 자의 바닥을 같이 잰다}

\begin{quote}
\textbf{조항} --- 새 채점 자를 만들면 \textbf{그 자에서 아무 손해도
없어야 할 대조}(여기서는 라벨 퍼짐을 맞춘 무작위 무리)를 같이 재서
\emph{자 자체가 깎는 양}을 적는다. 안 적으면 그 깎임을 모형 탓으로
읽는다.
\end{quote}

\section*{남은 것}

\begin{itemize}
\item \textbf{만화가 왜 반대인가} --- 갈래 안에서 더 잘하는 유일한
도메인이고 퍼짐을 맞추니 더 두드러진다.
\item \textbf{남은 $0.0266$ 중 펀딩을 빼면} 얼마인가 --- 펀딩은 노트 471
이 이미 설명했으니, 나머지 다섯만의 값이 진짜 남은 물음이다.
\item \textbf{노트 480$\cdot$481 의 축 수 격차도 같은 검사가 필요하다}
--- 축 많은/적은 두 무리의 라벨 퍼짐이 다를 수 있다.
\end{itemize}

\end{document}
